\subsection*{Materiales}
\begin{itemize}
\item Circuito desconocido con parámetros nominales $V_{eq}=12\unit{\volt}$,
$R_{eq}=1\unit{\kilo\ohm}$.
\item Amperimetro $\pm0,001\unit{\milli\ampere}$ y Voltimetro $\pm0,001\unit{\volt}$.
\item Resistencia variable.
\end{itemize}

\subsection*{Procedimiento}
La práctica se realizó con la herramienta digital {\it CircuitJS} por {\it Paul Falstad}
\cite{falstad}. Los materiales listados son equivalentes a los utilizados en el simulador y
se eligieron incertidumbres razonables para los elementos relevantes.

Se realizaron los tres circuitos mostrados en las Fig. \ref{fig:c1}---\ref{fig:c3}.
Para cada circuito se aumentó la resistencia variable $R_z$ entre \qtylist{1;8}{\kilo\ohm}
en intervalos de \qty{1}{\kilo\ohm}, y en cada intervalo se registró el amperaje y
voltaje marcado por los respectivos instrumentos.